\documentclass[12pt]{article}
\usepackage[utf8]{inputenc}
\usepackage{geometry}
\geometry{letterpaper, margin=1in}
\usepackage{amsmath}
\usepackage{graphicx} % For embedding images (Q1b)
\usepackage{enumitem}
\usepackage{booktabs} % For nice tables (Q3)
\usepackage{tikz}

\title{CS 412 Assignment 3 Report}
\author{Donggyu Park}
\date{\today}

\begin{document}

\maketitle

% ==========================================
% Question 1
% ==========================================
\section*{1. Decision Tree (25 points)}

\begin{enumerate}[label=(\alph*)]
    \item \textbf{(11 points) Compute the Information Gain for all of the features. Show all intermediate computations.}
    \\ The target is to classify "Spoiled".
    \\ \textbf{Entropy for "Spoiled":}

        $$ H(S) = -\sum_{i=1}^{n} p_i \log_2 p_i $$ 
  
        $$ H(S) = - \left(\frac{5}{11}\right) \log_2\left(\frac{5}{11}\right) - \left(\frac{6}{11}\right) \log_2\left(\frac{6}{11}\right) $$ 

        $$ H(S) = 0.9940302114769565 $$

    \textbf{Information Gain:} $IG(S, A) = H(S) - \sum_{v \in Values(A)} \frac{|S_v|}{|S|} H(S_v)$
    \begin{enumerate}[label=\roman*.]
        \item \textbf{Color} 
            $$H(\text{Red}) = - \left(\frac{2}{5}\right) \log_2\left(\frac{2}{5}\right) - \left(\frac{3}{5}\right) \log_2\left(\frac{3}{5}\right) = 0.9709505944546686$$
            $$H(\text{Yellow}) = - \left(\frac{3}{6}\right) \log_2\left(\frac{3}{6}\right) - \left(\frac{3}{6}\right) \log_2\left(\frac{3}{6}\right) = 1.0$$

            \textbf{Weighted average entropy for Color:}
            $$H(S_{\text{Color}}) = \frac{5}{11} \cdot 0.9709505944546686 + \frac{6}{11} \cdot 1.0 = 0.986795724752122$$

            \textbf{Information Gain for Color:}
            $$IG(S, \text{Color}) = H(S) - H(S_{\text{Color}}) = 0.9940 - 0.9868 = 0.00723448672483451$$

        \item \textbf{Texture}
            $$H(\text{Smooth}) = - \left(\frac{3}{5}\right) \log_2\left(\frac{3}{5}\right) - \left(\frac{2}{5}\right) \log_2\left(\frac{2}{5}\right) = 0.9709505944546686$$
            $$H(\text{Rough}) = - \left(\frac{2}{4}\right) \log_2\left(\frac{2}{4}\right) - \left(\frac{2}{4}\right) \log_2\left(\frac{2}{4}\right) = 1.0$$
            $$H(\text{Wrinkled}) = 0$$
            $$H(\text{Soft}) = 0$$

            \textbf{Weighted average entropy for Texture:}
            $$H(S_{\text{Texture}}) = \frac{5}{11} \cdot 0.97095 + \frac{4}{11} \cdot 1.0 + \frac{1}{11} \cdot 0 + \frac{1}{11} \cdot 0 = 0.8049775429339403$$

            \textbf{Information Gain for Texture:}
            $$IG(S, \text{Texture}) = H(S) - H(S_{\text{Texture}}) = 0.9940 - 0.8050 = 0.18905266854301617$$

        \item \textbf{Smell}
            $$H(1) = 0$$
            $$H(2) = - \left(\frac{2}{5}\right) \log_2\left(\frac{2}{5}\right) - \left(\frac{3}{5}\right) \log_2\left(\frac{3}{5}\right) = 0.9709505944546686$$
            $$H(3) = 0$$

            \textbf{Weighted average entropy for Smell:}
            $$H(S_{\text{Smell}}) = \frac{3}{11} \cdot 0 + \frac{5}{11} \cdot 0.97095 + \frac{3}{11} \cdot 0 = 0.4413411792975766$$

            \textbf{Information Gain for Smell:}
            $$IG(S, \text{Smell}) = H(S) - H(S_{\text{Smell}}) = 0.9940 - 0.4413 = 0.5526890321793798$$
    \end{enumerate}
    % TODO: Show Entropy(S) and Information Gain for Color, Texture, and Smell.
    
    \item \textbf{(5 points) Draw the full decision tree that would be learned for this dataset.}
    \\
\begin{center}
\begin{tikzpicture}[
  every node/.style = {shape=rectangle, rounded corners, draw, align=center, minimum width=1.2cm, minimum height=0.6cm},
  edge from parent/.style = {draw, -latex},
  level 1/.style={sibling distance=5cm, level distance=4cm},
  level 2/.style={sibling distance=5cm, level distance=3.5cm},
  level 3/.style={sibling distance=2cm}
]

  \node {\textbf{Smell}}
    child { node {Yes\\Yes\\Yes} 
      edge from parent node[left, draw=none] {1~~} 
    }
    child { node {\textbf{Texture}}
      child { node {Yes} 
        edge from parent node[left, draw=none] {Rough~~} 
      }
      child { node {No} 
        edge from parent node[left, draw=none] {Wrinkled~~} 
      }
      child { node {No} 
        edge from parent node[right, draw=none] {~~Soft} 
      }
      child { node {\textbf{Color}}
        child { node {Yes} 
          edge from parent node[left, draw=none] {Red~~} 
        }
        child { node {No} 
          edge from parent node[right, draw=none] {~~Yellow} 
        }
        edge from parent node[right, draw=none] {~~Smooth} 
      }
      edge from parent node[fill=white, draw=none] {2} 
    }
    child { node {No\\No\\No} 
      edge from parent node[right, draw=none] {~~3} 
    };

\end{tikzpicture}
\end{center}
    
    \item \textbf{(9 points) Fruit ID Feature Analysis}
    \begin{enumerate}[label=\roman*.]
        \item \textbf{What would the Information Gain for the Fruit ID feature be?}
        \\
        % TODO: Answer here.
        Fruit\_id will be a unique attribute.
        Thus if we split with Fruit\_ID there will be only one data point in each subset.
        The entropy of each subset will thus be 0,
        making the information gain be the same as the initial entropy of the dataset.
        
        \item \textbf{Would the standard decision tree algorithm (ID3) choose Fruit ID as the root?}
        \\
        % TODO: Answer here.
        Yes because this way, the information gain is maximized
        
        \item \textbf{Is Fruit ID a useful feature? Describe the fundamental flaw or bias in the Information Gain metric.}
        \\
        % TODO: Explain why IG favors attributes with many distinct values.
        No. If an unseen fruit\_id is given, the tree won't know how to appropriately classify it, 
        as the branches were constructed based on fruit\_1 to fruit\_11. 
        We can conclude that Information Gain is a greedy algorithm that is biased to the features with more unique values.


    \end{enumerate}
\end{enumerate}

\newpage
% ==========================================
% Question 2
% ==========================================
\section*{2. Linear Models (30 points)}

\textbf{Part A}
\begin{enumerate}[label=(\alph*)]
    \item \textbf{(5 pts) Why does the decision boundary in Figure B provide poorer classification performance than the boundary in Figure A?}
    \\
    % TODO: Answer here.
    If we only look at the data points in the plots, the decision boundary in Figure B does not provide poor classification performance than A.
    Boundary A fails to classify four points of squares, while Boundary B fials to classify three points of square.
    However, the decision boundary in Figure B is influenced by an extreme point at (7,-6) which is an outlier.
    This outlier point distorts the decision boundary, making it less representative of the overall data distribution.
    Due to this distortion the generalization performance is degraded.

    \item \textbf{(5 pts) Explain how the square loss function contributes to the issue observed in Figure B.}
    \\
    % TODO: Answer here.
    We can see that Boundary B is influenced by the outlier point at (7,-6) with its slope smaller than that of Boundary A.
    This is because the square loss function penalizes larger errors that are far from the boundary more heavily than smaller ones.

    
    \item \textbf{(5 pts) Suppose the extreme point in the right plot were removed. Describe how the learned decision boundary would likely change and explain why.}
    \\
    % TODO: Answer here.
    Since the outlier point at (7,-6) is removed, 
    the decision boundary would likely shift to better fit the remaining data points like Boundary A.
\end{enumerate}

\vspace{1em}
\textbf{Part B}
\begin{enumerate}[label=(\alph*), start=4]
    \item \textbf{(5 pts) Describe how you might use an SVM as a classifier to predict whether an item belongs to the +1 or -1 class in Figure 1.}
    \\
    % TODO: Explain the use of a non-linear kernel (e.g., RBF or Polynomial) for the circular boundary.
    Data points in Figure 1 are distrubuted in a circular pattern, making it impossible to separate the with a linear boundary.
    Thus, a non linear kernel such as RBF or polynomial must be introduced to transform the data into a high dimesion space.
    Then, a linear hyperplane can be used to separate the data in the space.

    \item \textbf{(5 pts) In SVM, it is often the case that data points cannot be well-separated via a "hard margin". Can you provide a solution to solve this problem?}
    \\
    % TODO: Explain soft-margin SVM and the introduction of slack variables.
    If the data is noisy, hard margin may fail because it tries to strictly separate the "noisy" data.
    Therefore, by introducing soft margin, we have to allow some misclassifications to form a less-sensitive decision boundary.
    This is achieved by introducing slack variables that allow some data points to be on the wrong side.

    \item \textbf{(5 pts) In a soft-margin SVM, the parameter C controls the trade-off...}
    \begin{itemize}
        \item \textbf{What happens when C is very large?}
        \\
        % TODO: Answer here.
        Big penalties are imposed on misclassifications. 
        The decision boundary is likely to be hard margin, which may lead to overfitting if the data is noisy.
        
        \item \textbf{What happens when C is very small?}
        \\
        % TODO: Answer here.
        Little penalty is imposed on misclassifications.
        The decision boundary is likely to be very flexible, which may lead to underfitting.

    \end{itemize}
\end{enumerate}

\newpage
% ==========================================
% Question 3
% ==========================================
\section*{3. Naive Bayes Classifier (25 points)}

\begin{enumerate}[label=(\alph*)]
    \item \textbf{(20 points) Compute the likelihood tables for all three features.}
    \\
    % TODO: Fill in the probabilities based on the dataset.
    
    \textbf{Table 2: Likelihood Table for Nausea}
    \begin{center}
    \begin{tabular}{lcc}
        \toprule
        & \multicolumn{2}{c}{\textbf{Has Food Poisoning}} \\
        \cmidrule(lr){2-3}
        \textbf{Nausea} & \textbf{Positive} & \textbf{Negative} \\
        \midrule
        \textbf{Yes} & $\frac{3}{4}$ & $\frac{1}{3}$ \\
        \textbf{No}  & $\frac{1}{4}$ & $\frac{2}{3}$ \\
        \bottomrule
    \end{tabular}
    \end{center}

    \vspace{1em}
    \textbf{Table 3: Likelihood Table for Fever}
    \begin{center}
    \begin{tabular}{lcc}
        \toprule
        & \multicolumn{2}{c}{\textbf{Has Food Poisoning}} \\
        \cmidrule(lr){2-3}
        \textbf{Fever} & \textbf{Positive} & \textbf{Negative} \\
        \midrule
        \textbf{Yes} & $\frac{3}{4}$ & $\frac{1}{3}$ \\
        \textbf{No}  & $\frac{1}{4}$ & $\frac{2}{3}$ \\
        \bottomrule
    \end{tabular}
    \end{center}

    \vspace{1em}
    \textbf{Table 4: Likelihood Table for Stomach Cramps}
    \begin{center}
    \begin{tabular}{lcc}
        \toprule
        & \multicolumn{2}{c}{\textbf{Has Food Poisoning}} \\
        \cmidrule(lr){2-3}
        \textbf{Stomach Cramps} & \textbf{Positive} & \textbf{Negative} \\
        \midrule
        \textbf{Yes} & $\frac{3}{4}$ & $\frac{0}{3}$ \\
        \textbf{No}  & $\frac{1}{4}$ & $\frac{3}{3}$ \\
        \bottomrule
    \end{tabular}
    \end{center}

    \item \textbf{(5 points) Classify: Nausea = No, Fever = Yes, and Stomach Cramps = No. Show intermediate calculations.}
    \\ \\
    % TODO: Show P(Positive|X) and P(Negative|X) calculations using Bayes theorem, then compare.
    \textbf{Prior Probabilities} \\
    $P(\text{Positive}) = 4/7$ \\
    $P(\text{Negative}) = 3/7$ \\ \\

    \textbf{Likelihoods} \\
    \textbf{Positive Class:} \\
    $P(\text{Nausea} = \text{No}|\text{Positive}) = \frac{1}{4}$ \\
    $P(\text{Fever} = \text{Yes}|\text{Positive}) = \frac{3}{4}$ \\
    $P(\text{Stomach Cramps} = \text{No}|\text{Positive}) = \frac{1}{4}$ \\ \\
    \textbf{Negative Class:} \\
    $P(\text{Nausea} = \text{No}|\text{Negative}) = \frac{2}{3}$ \\
    $P(\text{Fever} = \text{Yes}|\text{Negative}) = \frac{1}{3}$ \\
    $P(\text{Stomach Cramps} = \text{No}|\text{Negative}) = \frac{3}{3}$ \\ \\

    \textbf{Score(Positive)} \\
    $$\text{Score(Pos)} = \frac{4}{7} \times \frac{1}{4} \times \frac{3}{4} \times \frac{1}{4} = \frac{3}{112} = 0.0268$$
    \textbf{Score(Negative)} \\
    $$\text{Score(Neg)} = \frac{3}{7} \times \frac{2}{3} \times \frac{1}{3} \times \frac{3}{3} = \frac{6}{63} = 0.0952$$

    \textbf{Normalization} \\
    $P(\text{Positive} \mid X) = \frac{3/112}{41/336} = \frac{9}{41} = \textbf{0.2195}$\\
    $P(\text{Negative} \mid X) = \frac{2/21}{41/336} = \frac{32}{41} = \textbf{0.7805}$


\end{enumerate}

\newpage
% ==========================================
% Question 4
% ==========================================
\section*{4. Posterior Probabilities (20 points)}

\begin{enumerate}[label=(\alph*)]
    \item \textbf{(8 points) Draw a candy (Candy1) = "lime". What are the posterior probabilities of each type of bag?}
    \\
    % TODO: Show Bayes rule calculations for p(h1|lime), p(h2|lime), p(h3|lime).
    \textbf{Bayes Rule:}\\
    $$P(h_i | \text{lime}_1) = \frac{P(\text{lime}_1 | h_i) P(h_i)}{P(\text{lime}_1)}$$

    \textbf{Prior Probabilities} \\
    $P(h_1) = 1/4$ \\
    $P(h_2) = 1/2$ \\
    $P(h_3) = 1/4$

    \textbf{Likelihoods} \\
    $P(\text{lime}_1 | h_1) = 0$ \\
    $P(\text{lime}_1 | h_2) = 50/100 = 1/2$ \\
    $P(\text{lime}_1 | h_3) = 100/100 = 1$

    \textbf{Marginal Probability (Law of Total Probability)} \\
    $$P(\text{lime}_1) = (0 \times 1/4) + (1/2 \times 1/2) + (1 \times 1/4) = 0 + 1/4 + 1/4 = 1/2$$

    \textbf{Posterior Probabilities} \\
    $P(h_1 | \text{lime}_1) = \frac{0 \times 1/4}{1/2} = \textbf{0}$ \\
    $P(h_2 | \text{lime}_1) = \frac{(1/2) \times (1/2)}{1/2} = \textbf{1/2}$ \\
    $P(h_3 | \text{lime}_1) = \frac{1 \times 1/4}{1/2} = \textbf{1/2}$



    \item \textbf{(8 points) Draw another candy (Candy2) = "lime" (without replacement). What are the new posterior probabilities?}
    \\
    % TODO: Show Bayes rule calculations considering the updated bag contents.

    \textbf{New Priors (from part a):} \\
    $P'(h_1) = 0$ \\
    $P'(h_2) = 1/2$ \\
    $P'(h_3) = 1/2$

    \textbf{New Likelihoods ($Candy_2 = \text{lime}$ given $Candy_1 = \text{lime}$):} \\
    Since one lime candy was removed, the bags now have 99 candies.\\
    $P(\text{lime}_2 | h_1, \text{lime}_1) = 0$ \\
    $P(\text{lime}_2 | h_2, \text{lime}_1) = 49/99$ \\
    $P(\text{lime}_2 | h_3, \text{lime}_1) = 99/99 = 1$

    \textbf{Marginal Probability for the second draw:} \\
    $$P(\text{lime}_2 | \text{lime}_1) = (0 \times 0) + \left(\frac{49}{99} \times \frac{1}{2}\right) + \left(1 \times \frac{1}{2}\right)$$
    $$P(\text{lime}_2 | \text{lime}_1) = \frac{49}{198} + \frac{99}{198} = \frac{148}{198} = \frac{74}{99}$$

    \textbf{Posterior Probabilities:} \\
    $P(h_1 | \text{lime}_1, \text{lime}_2) = 0$ \\
    $P(h_2 | \text{lime}_1, \text{lime}_2) = \frac{(49/99) \times (1/2)}{74/99} = \frac{49/198}{148/198} = \frac{49}{148}$ \\
    $P(h_3 | \text{lime}_1, \text{lime}_2) = \frac{1 \times (1/2)}{74/99} = \frac{99/198}{148/198} = \frac{99}{148}$
        
    \item \textbf{(4 points) Are your computations in (b) above using conditional independence in any form?}
    \\
    % TODO: Answer Yes/No and justify based on the "without replacement" condition.
    No, the computations in (b) do not use conditional independence. 
    The outcome of the first draw physically alters the contents of the bag.
\end{enumerate}

\end{document}